\documentclass[tikz,border=8pt]{standalone}
\usepackage{ctex}
\usepackage{amsmath}
\usetikzlibrary{calc, arrows.meta}
\begin{document}
\begin{tikzpicture}[scale=1.5, >={Stealth[length=2mm]}]
  % 坐标轴
  \draw[->] (-2.1,0) -- (2.2,0) node[right,font=\small] {$x$};
  \draw[->] (0,-0.35) -- (0,2.15) node[above,font=\small] {$y$};
  \foreach \x in {-2,-1,1,2} {\draw (\x,-0.04) -- (\x,0.04) node[below=2pt,font=\scriptsize] {$\x$};}

  % 支撑直线（斜率取遍 [-1,1] 的若干条），全部从下方托住 |x|，相交于原点
  \foreach \s/\c in {-1/red!55, -0.5/orange!70!black, 0/green!55!black, 0.5/teal, 1/blue!70} {
    \draw[\c, thin] (-1.75, {-1.75*\s}) -- (1.75, {1.75*\s});
  }

  % |x| 主图形（黑色粗线）
  \draw[black, very thick] (-2,2) -- (0,0) -- (2,2);

  % 尖点
  \fill[black] (0,0) circle (1.3pt);

  % 极端支撑线标注（斜率 ±1 与左右两支重合）
  \node[blue!70, font=\scriptsize, anchor=south west] at (1.15,1.05) {斜率 $+1$};
  \node[red!55, font=\scriptsize, anchor=south east] at (-1.15,1.05) {斜率 $-1$};
  \node[green!55!black, font=\scriptsize, anchor=south] at (1.55,0.02) {斜率 $0$};

  % 标题
  \node[font=\bfseries] at (0.1,2.55) {次梯度：$|x|$ 在尖点 $x=0$ 处的支撑直线（斜率取遍 $[-1,1]$）};

  % 半透明注释框，置于左下空白区，避开图形
  \node[font=\scriptsize, fill=yellow!18, draw=yellow!60!black, rounded corners=2pt,
        align=center, fill opacity=0.9, text opacity=1, anchor=north]
    at (-1.05,-0.12)
    {$\partial|x|\big|_{x=0}=[-1,1]$\\[1pt] 每条斜率 $\in[-1,1]$ 的直线\\都从下方托住 $|x|$};

\end{tikzpicture}
\end{document}
